\documentclass[11pt,a4paper]{article}
\usepackage[spanish,es-noquoting]{babel}
\usepackage[top=2.5cm, bottom=2.5cm, left=2.5cm, right=2.5cm]{geometry}
\usepackage{fontspec}
\usepackage{amsmath}
\usepackage[version=4]{mhchem}
\usepackage{enumitem}
\usepackage{microtype}

\setmainfont{Calibri}
\linespread{1.15}

\begin{document}
\pagestyle{empty}

\noindent\textbf{1. Resolución del ejercicio de neutralización (\ce{H2SO4 + NaOH}):}

\vspace{0.2cm}
\noindent\textbf{Cálculo de masas molares ($M$):}
\begin{itemize}[leftmargin=1.5em, itemsep=2pt]
    \item $M(\ce{H2SO4}) = 2 \cdot 1 + 32 + 4 \cdot 16 = 98\text{ g/mol}$
    \item $M(\ce{NaOH}) = 23 + 16 + 1 = 40\text{ g/mol}$
    \item $M(\ce{Na2SO4}) = 2 \cdot 23 + 32 + 4 \cdot 16 = 142\text{ g/mol}$
    \item $M(\ce{H2O}) = 2 \cdot 1 + 16 = 18\text{ g/mol}$
\end{itemize}

\vspace{0.2cm}
\noindent\textbf{a) Ajuste de la ecuación química:}
\[
\ce{H2SO4 (aq) + 2 NaOH (aq) -> Na2SO4 (aq) + 2 H2O (l)}
\]

\vspace{0.1cm}
\noindent\textbf{b) Determinación del reactivo limitante y reactivo en exceso:}
\begin{itemize}[leftmargin=1.5em, itemsep=3pt]
    \item Moles iniciales de cada reactivo:
    \[
    n(\ce{H2SO4}) = \frac{49\text{ g}}{98\text{ g/mol}} = 0{,}50\text{ mol de \ce{H2SO4}}
    \]
    \[
    n(\ce{NaOH}) = \frac{32\text{ g}}{40\text{ g/mol}} = 0{,}80\text{ mol de \ce{NaOH}}
    \]
    \item Relación estequiométrica: $1\text{ mol de \ce{H2SO4}}$ reacciona con $2\text{ mol de \ce{NaOH}}$.
    \item Cantidad de \ce{H2SO4} necesaria para consumir todo el \ce{NaOH}:
    \[
    n(\ce{H2SO4})_{\text{necesario}} = 0{,}80\text{ mol \ce{NaOH}} \cdot \frac{1\text{ mol \ce{H2SO4}}}{2\text{ mol \ce{NaOH}}} = 0{,}40\text{ mol de \ce{H2SO4}}
    \]
    \item Como disponemos de $0{,}50\text{ mol}$ de \ce{H2SO4} ($0{,}50 > 0{,}40$), el \textbf{reactivo limitante} es el hidróxido de sodio (\ce{NaOH}) y el \textbf{reactivo en exceso} es el ácido sulfúrico (\ce{H2SO4}).
\end{itemize}

\vspace{0.1cm}
\noindent\textbf{c) Masa de sulfato de sodio (\ce{Na2SO4}) formada:}
\begin{itemize}[leftmargin=1.5em, itemsep=3pt]
    \item Calculamos los moles y la masa a partir del reactivo limitante ($0{,}80\text{ mol de \ce{NaOH}}$):
    \[
    n(\ce{Na2SO4}) = 0{,}80\text{ mol \ce{NaOH}} \cdot \frac{1\text{ mol \ce{Na2SO4}}}{2\text{ mol \ce{NaOH}}} = 0{,}40\text{ mol de \ce{Na2SO4}}
    \]
    \[
    m(\ce{Na2SO4}) = 0{,}40\text{ mol} \cdot 142\text{ g/mol} = \mathbf{56{,}8\text{ g de \ce{Na2SO4}}}
    \]
\end{itemize}

\vspace{0.1cm}
\noindent\textbf{d) Cantidad y masa del reactivo en exceso sobrante:}
\begin{itemize}[leftmargin=1.5em, itemsep=3pt]
    \item Moles de \ce{H2SO4} que reaccionan: $0{,}40\text{ mol}$.
    \item Moles sobrantes: $n_{\text{sobrante}} = 0{,}50\text{ mol} - 0{,}40\text{ mol} = \mathbf{0{,}10\text{ mol de \ce{H2SO4}}}$.
    \item Masa sobrante: $m_{\text{sobrante}} = 0{,}10\text{ mol} \cdot 98\text{ g/mol} = \mathbf{9{,}8\text{ g de \ce{H2SO4}}}$.
\end{itemize}

\vspace{0.1cm}
\noindent\textbf{e) Masa de agua formada y comprobación de la ley de Lavoisier:}
\begin{itemize}[leftmargin=1.5em, itemsep=3pt]
    \item Moles de agua producidos: $0{,}80\text{ mol \ce{NaOH}} \cdot \frac{2\text{ mol \ce{H2O}}}{2\text{ mol \ce{NaOH}}} = 0{,}80\text{ mol de \ce{H2O}}$.
    \item Masa de agua: $m(\ce{H2O}) = 0{,}80\text{ mol} \cdot 18\text{ g/mol} = \mathbf{14{,}4\text{ g de \ce{H2O}}}$.
    \item Comprobación de la conservación de la masa:
    \[
    m_{\text{inicial}} = 49{,}0\text{ g} + 32{,}0\text{ g} = 81{,}0\text{ g}
    \]
    \[
    m_{\text{final}} = 56{,}8\text{ g (\ce{Na2SO4})} + 14{,}4\text{ g (\ce{H2O})} + 9{,}8\text{ g (\ce{H2SO4} sobrante)} = 81{,}0\text{ g}
    \]
    Se cumple la ley de Lavoisier ($m_{\text{inicial}} = m_{\text{final}} = 81{,}0\text{ g}$).
\end{itemize}

\end{document}
